% TOP_PROBLEM: {"id": 30006785, "title": "Polynomial Freiman-Ruzsa Conjecture for General Abelian Groups", "category_id": 2, "category": {"id": 2, "name": "combinatorics", "display_name": "Combinatorics", "description": "Counting problems, graph theory, discrete structures.", "slug": "combinatorics", "order_index": 2, "created_at": "2026-07-31T15:26:25.671Z"}, "status": "open"}
% FIRST_POSTED: 2026-09-27
% !TeX program = lualatex
% Compile twice: lualatex open_problems_500.tex
% All references and prose are embedded; no BibTeX database or external inputs.
\documentclass[11pt,a4paper]{article}
\usepackage[margin=24mm,headheight=14pt]{geometry}
\usepackage{fontspec}
\setmainfont{Latin Modern Roman}
\setsansfont{Latin Modern Sans}
\setmonofont{Latin Modern Mono}
\usepackage{amsmath,amssymb,mathtools}
\usepackage{unicode-math}
\setmathfont{Latin Modern Math}
\usepackage{microtype}
\usepackage{enumitem,xurl,needspace}
\usepackage[dvipsnames]{xcolor}
\usepackage{hyperref}
\hypersetup{unicode=true,colorlinks=true,linkcolor=MidnightBlue,urlcolor=MidnightBlue,
 pdftitle={500 Mathematical Problem Statements},pdfauthor={Source-annotated compilation},bookmarksnumbered=true}
\usepackage{bookmark}
\usepackage{tocloft}
\setlength{\cftsecnumwidth}{3em}
\cftsetpnumwidth{2.5em}
\cftsetrmarg{4em}
\usepackage{titlesec}
\titleformat{\section}{\Large\bfseries\raggedright}{\thesection}{0.7em}{}
\usepackage{fancyhdr}
\pagestyle{fancy}\fancyhf{}
\fancyhead[L]{\small\sffamily Mathematical problem statements}
\fancyhead[R]{\small Source edition 22}
\fancyfoot[C]{\thepage}
\setlength{\parindent}{0pt}
\setlength{\parskip}{5pt plus 1pt}
\setlength{\emergencystretch}{3em}
\setcounter{tocdepth}{1}
\setcounter{secnumdepth}{1}
\setlist[itemize]{leftmargin=3em,itemsep=5pt,topsep=4pt,parsep=0pt}
\allowdisplaybreaks
\newcommand{\N}{\mathbb{N}}
\newcommand{\Z}{\mathbb{Z}}
\newcommand{\Q}{\mathbb{Q}}
\newcommand{\R}{\mathbb{R}}
\newcommand{\C}{\mathbb{C}}
\newcommand{\F}{\mathbb{F}}
\newcommand{\A}{\mathbb{A}}
\newcommand{\PP}{\mathbb P}
\newcommand{\E}{\mathbb{E}}
\newcommand{\eps}{\varepsilon}
\newcommand{\dd}{\,\mathrm d}
\newcommand{\sref}[2]{\hyperlink{source:#1:#2}{[S#2]}}
\newcommand{\eref}[2]{\hyperlink{extra:#1:#2}{[E#2]}}
% Turn the next switch off for a shorter reading copy. The quotations preserve
% every exactTarget field, including qualifications omitted from a short summary.
\newif\ifshowcatalogue
\showcataloguetrue
\newcommand{\cataloguescope}[1]{\ifshowcatalogue
 \paragraph{Exact scope in the supplied catalogue.}
 \begin{quote}\small #1\end{quote}\fi}

% Compile twice with: lualatex 30006785.tex
\numberwithin{equation}{section}
\hypersetup{pdftitle={Research attempt -- problem 30006785},pdfauthor={Research notebook; original compilation retained}}
\fancyhead[L]{\small\sffamily Research notebook}
\fancyhead[R]{\small\sffamily Problem 30006785 | 27 September 2026}
\begin{document}
\setcounter{section}{0}
% ================================================================
% problemId: 30006785
% Canonical problem ID: 30006785
% exactTarget SHA-256: 31d392b5ed5cba6b34cf57e54bce24d30dc37303907d771c70a01dfe023c14d1
\clearpage
\section*{\texorpdfstring{Polynomial Freiman-Ruzsa Conjecture for General Abelian Groups}{Polynomial Freiman-Ruzsa Conjecture for General Abelian Groups}}\label{30006785}
\subsection{Definitions and mathematical statement}
For a nonempty finite set $A$ in an abelian group, define $A+A=\{a+a':a,a'\in A\}$ and suppose $|A+A|\le K|A|$, $K\ge1$. A centered convex coset progression of dimension $d$ has the form
$P=H+\phi(C\cap\Z^d)$, where $H$ is a finite subgroup, $\phi:\Z^d\to G$ is a homomorphism, and $C\subset\R^d$ is a centrally symmetric convex body centered at zero. The assertion is that an absolute $C_0$ allows
\[
 A\subseteq\bigcup_{t\in T}(t+P),\quad |T|\le K^{C_0},
 \quad d\le C_0\log K,\quad |P|\le K^{C_0}|A|.
\]
No bounded exponent or torsion-free restriction is placed on the ambient group, and the selected progression is convex rather than necessarily rectangular.
\par\smallskip\textit{Formulation and concept references:} \sref{30006785}{1}.
\cataloguescope{Prove that there is an absolute constant C such that for every finite subset A of an arbitrary abelian group with |A+A| <= K|A|, the set A can be covered by exp(C log K) translates of a centred convex coset progression of dimension at most C log K and size at most exp(C log K)|A|.}
\subsection{Short English statement}
Can every finite set with small additive doubling in any abelian group be covered efficiently by translates of a low-dimensional convex coset progression?
% BEGIN RESEARCH ADDITION ATTEMPT-195
\subsection{Research attempt}

\paragraph{Current frontier.}
The general-group statement concerns convex coset progressions with absolute polynomial bounds, as in Sanders' formulation \sref{30006785}{1}. The bounded-torsion breakthrough proves polynomial covering for groups of a fixed exponent $m$, with a bound of the form $(2K)^{O(m^3)}$ \eref{30006785}{1}, Theorem~1.1. The dependence on $m$ prevents a direct reduction of arbitrary groups by a large-modulus quotient. We investigate a geometric large-subset criterion and prove its equivalence, up to absolute polynomial losses, to the requested covering statement.

\paragraph{Attack route.}
An entropy argument might find a large structured piece without constructing a cover immediately. A covering lemma can then transfer it to all of $A$, provided enlargement of its convex progression has controlled size. Alternatively, one might reduce a finitely generated ambient group modulo a large exponent and invoke the bounded-torsion theorem. The first route reduces the target to a precise large-piece statement; the second fails to retain an absolute exponent. We check both quantitatively, including the endpoint $K=1$.

\paragraph{Attempt.}
Assume first $K>1$. Suppose there are absolute nonnegative constants $a,b,c$ and, for every set $A$ in the question, a subset $B\subseteq A$ and a translate of a centered convex coset progression
\[
 B\subseteq x+P,\qquad P=H+\phi(C\cap\Z^d)
\]
such that
\begin{equation}\label{eq195:chunk}
 |B|\ge K^{-a}|A|,\qquad |P|\le K^b|A|,
 \qquad d\le c\log K.
\end{equation}
This is the missing large-piece assertion, not an assumption known for all groups.

Choose a maximal set $T\subseteq A$ for which the translates $t+B$, $t\in T$, are pairwise disjoint. They are contained in $A+B\subseteq A+A$, so
\[
 |T||B|\le K|A|,\qquad |T|\le K^{a+1}.
\]
Maximality gives $A\subseteq T+B-B$: otherwise $a+B$ could be added. Also
\[
 B-B\subseteq P-P\subseteq P_2:=H+\phi(2C\cap\Z^d).
\]
To finish, the size of $P_2$ must be controlled without assuming $\phi$ injective or the progression proper.

\textbf{Convex doubling lemma.} For a centered symmetric convex body $C\subset\R^d$,
\begin{equation}\label{eq195:convex}
 2C\cap\Z^d\subseteq Z+(C\cap\Z^d),\qquad |Z|\le5^d
\end{equation}
for a finite $Z\subseteq2C\cap\Z^d$.
Choose $Z$ maximal with $z-z'\notin C$ for distinct selected points. The sets $z+\tfrac12\operatorname{int}C$ are disjoint: overlap would put the difference in $\operatorname{int}C$. They are all contained in $\tfrac52 C$. Comparing volumes gives $|Z|\le5^d$. Maximality gives the covering, and the differences are integral. The case $d=0$ is immediate. No statement about the number of lattice points by volume is used, so thin bodies cause no error.

Mapping \eqref{eq195:convex} by $\phi$ and adding $H$ gives
\[
 |P_2|\le5^d|P|\le K^{b+c\log5}|A|.
\]
Thus $P_2$ is still a centered convex coset progression of dimension $d$, and at most $K^{a+1}$ translates cover $A$. Enlarging a single absolute constant to dominate $a+1,b+c\log5,c$ gives exactly the polynomial and logarithmic form in the target.

Conversely, a cover by at most $K^C$ translates of a progression of size at most $K^C|A|$ has one translate meeting $A$ in at least $K^{-C}|A|$ points. Taking that intersection as $B$ proves \eqref{eq195:chunk}. Hence large-piece control and covering control are equivalent up to explicit absolute exponent losses. This is an independently written proof of a standard covering mechanism, not a claim that the missing large piece has been constructed.

\textbf{Endpoint check.} If $|A+A|=|A|$, translate so $0\in B=A-a$. Then $B\subseteq B+B$ and equal cardinality forces $B+B=B$. A finite set containing zero and closed under addition is a subgroup: for any $b$, repetition among its nonnegative multiples gives a positive multiple zero and hence its inverse in $B$. Therefore $A$ is one coset of a finite subgroup. The target holds with dimension zero, one translate, and $|P|=|A|$, respecting the exact $K=1$ bounds rather than replacing $K$ by $2K$.

\textbf{Why a large finite quotient loses uniformity.} The subgroup generated by $A$ is finitely generated, so it has form $\Z^r\oplus T_0$. One can choose a modulus $q$ making reduction injective on the finite set $A+A$; the quotient has exponent dividing $\operatorname{lcm}(q,\exp T_0)$. Applying \eref{30006785}{1} then incurs an exponent depending on this number, not only on $K$. For the interval $A=\{0,\ldots,N-1\}\subset\Z$, $K=(2N-1)/N<2$, while injection on $A+A$ requires $q\ge2N-1$. Thus the theorem's $(2K)^{O(q^3)}$ loss is not bounded by $K^{C_0}$ with an absolute $C_0$. A small modulus avoids that exponent growth only by losing information that must then be recovered through a separate lifting theorem.

\paragraph{Stress test.}
The covering uses $A+B\subseteq A+A$, not an unjustified estimate for an unrelated set. The volume argument proves a lattice translate covering, so it remains valid after a noninjective homomorphism and in the presence of torsion. No conversion to rectangular progressions is made: such a conversion could introduce a dimension-dependent loss outside the stated target. A high affine rank of $A$ is not itself a counterexample, since the target permits polynomially many translates; sparse simplexes illustrate why a full-span argument would be too strong. The bounded-torsion result remains a genuine major special case but supplies no modulus-uniform version by the quotient trick above.

\paragraph{Outcome.}
\textbf{Outcome: Conditional reduction.}
It suffices to find a polynomially large subset of $A$ inside a polynomial-size convex coset progression of dimension $O(\log K)$. The passage from that subset to the exact requested cover is proved with explicit losses, including $K=1$. The unresolved input is the absolute large-piece theorem for arbitrary groups; the fixed-exponent theorem and large-quotient reduction do not establish it.

% END RESEARCH ADDITION ATTEMPT-195
\subsection{Sources}
\begin{itemize}\raggedright
\item[\textnormal{[S1]}]\hypertarget{source:30006785:1}{} Tom Sanders, "The structure theory of set addition revisited," Bulletin (New Series) of the American Mathematical Society, Volume 50, Number 1, January 2013, Pages 93-127; S 0273-0979(2012)01392-7, published electronically October 2, 2012.
\url{https://www.ams.org/journals/bull/2013-50-01/S0273-0979-2012-01392-7/S0273-0979-2012-01392-7.pdf}.
{\small\textit{Catalogue formulation source.}}
% BEGIN RESEARCH ADDITION REFERENCES-195
\item[\textnormal{[E1]}]\hypertarget{extra:30006785:1}{} W. T. Gowers, Ben Green, Freddie Manners and Terence Tao, \emph{Marton\textquotesingle s conjecture in abelian groups with bounded torsion}, arXiv:2404.02244v2, Theorem 1.1; explicit dependence on the ambient exponent.
\url{https://arxiv.org/html/2404.02244v2}.
{\small\textit{Research reference; consulted 27 September 2026.}}
% END RESEARCH ADDITION REFERENCES-195
\end{itemize}

\end{document}
